\section{支撑材料与复现说明}
\subsection{文件用途与运行方法}
支撑材料包为\texttt{cumcm2026A\_support.zip}。主要文件包括：\texttt{src/}下的物性、空气湿度、边界、有限体积求解及导出模块；\texttt{scripts/}下的求解、验证和图表生成入口；\texttt{tests/}下的核验程序；\texttt{results/result1.xlsx}至\texttt{result4.xlsx}、数值摘要和验证报告；\texttt{paper/}下的论文源文件、自动生成表格与矢量插图；\texttt{docs/model/}下的模型规格和文献依据。具体清单随支撑包附带，原始题目附件按目录\texttt{problem/}读取。

本次使用Python、NumPy、SciPy、pandas、openpyxl和Matplotlib完成求解与结果输出，以XeLaTeX编译中文论文。已配置环境时，在项目根目录执行：
\begin{lstlisting}[language={},numbers=none]
conda run -n 26A python scripts/run_pipeline.py
\end{lstlisting}
该入口依次生成主方案、附件拟合方案、验证报告、论文表图、源码附录、PDF与支撑包。单独计算四问可执行\texttt{python run\_all.py}；单独编译论文使用\texttt{latexmk -xelatex}，输出目录指定为\texttt{paper/build}。
所有数值先以SI单位计算，写出时再换算；论文表格从结果工作簿或同一次计算的连续终点记录读取，摘要和正文关键数字通过自动生成的宏引用。原始题目数据不修改，失败的积分和格式校验会中止计算。

\subsection{AI辅助使用说明}
本解题过程使用AI辅助核对题面与格式、检索文献、推导和审查模型、编写Python程序及组织LaTeX正文。论文中数值由本地程序实际运行得到，文献按可核验来源列示；AI生成的数值建议不直接作为结果。此次生成不包含参赛身份、学校、签名或任何提交操作。

\subsection{完整源程序}
以下源程序从项目实际文件直接载入；相同代码也以可运行的文本文件提供。论文排版中的自动换行不改变支撑包内的文件内容。
% Generated by paper/build_sources.py; source is included directly.
\subsubsection*{\texttt{run\_all.py}}
\lstinputlisting[language=Python]{../run_all.py}
\subsubsection*{\texttt{src/common/boundary.py}}
\lstinputlisting[language=Python]{../src/common/boundary.py}
\subsubsection*{\texttt{src/common/correlations.py}}
\lstinputlisting[language=Python]{../src/common/correlations.py}
\subsubsection*{\texttt{src/common/export.py}}
\lstinputlisting[language=Python]{../src/common/export.py}
\subsubsection*{\texttt{src/common/io\_excel.py}}
\lstinputlisting[language=Python]{../src/common/io_excel.py}
\subsubsection*{\texttt{src/common/numerical\_settings.py}}
\lstinputlisting[language=Python]{../src/common/numerical_settings.py}
\subsubsection*{\texttt{src/common/paths.py}}
\lstinputlisting[language=Python]{../src/common/paths.py}
\subsubsection*{\texttt{src/common/plotting.py}}
\lstinputlisting[language=Python]{../src/common/plotting.py}
\subsubsection*{\texttt{src/common/problem\_spec.py}}
\lstinputlisting[language=Python]{../src/common/problem_spec.py}
\subsubsection*{\texttt{src/common/psychrometrics.py}}
\lstinputlisting[language=Python]{../src/common/psychrometrics.py}
\subsubsection*{\texttt{src/common/radial.py}}
\lstinputlisting[language=Python]{../src/common/radial.py}
\subsubsection*{\texttt{src/common/solution.py}}
\lstinputlisting[language=Python]{../src/common/solution.py}
\subsubsection*{\texttt{src/common/units.py}}
\lstinputlisting[language=Python]{../src/common/units.py}
\subsubsection*{\texttt{src/q1/solver.py}}
\lstinputlisting[language=Python]{../src/q1/solver.py}
\subsubsection*{\texttt{src/q2/solver.py}}
\lstinputlisting[language=Python]{../src/q2/solver.py}
\subsubsection*{\texttt{src/q3/solver.py}}
\lstinputlisting[language=Python]{../src/q3/solver.py}
\subsubsection*{\texttt{src/q4/solver.py}}
\lstinputlisting[language=Python]{../src/q4/solver.py}
\subsubsection*{\texttt{scripts/\_bootstrap.py}}
\lstinputlisting[language=Python]{../scripts/_bootstrap.py}
\subsubsection*{\texttt{scripts/\_runner.py}}
\lstinputlisting[language=Python]{../scripts/_runner.py}
\subsubsection*{\texttt{scripts/audit\_q4\_boundary\_layer.py}}
\lstinputlisting[language=Python]{../scripts/audit_q4_boundary_layer.py}
\subsubsection*{\texttt{scripts/build\_paper\_assets.py}}
\lstinputlisting[language=Python]{../scripts/build_paper_assets.py}
\subsubsection*{\texttt{scripts/fit\_ambient\_temperature.py}}
\lstinputlisting[language=Python]{../scripts/fit_ambient_temperature.py}
\subsubsection*{\texttt{scripts/make\_figs.py}}
\lstinputlisting[language=Python]{../scripts/make_figs.py}
\subsubsection*{\texttt{scripts/make\_tables.py}}
\lstinputlisting[language=Python]{../scripts/make_tables.py}
\subsubsection*{\texttt{scripts/package\_delivery.py}}
\lstinputlisting[language=Python]{../scripts/package_delivery.py}
\subsubsection*{\texttt{scripts/parse\_attachments.py}}
\lstinputlisting[language=Python]{../scripts/parse_attachments.py}
\subsubsection*{\texttt{scripts/probe\_templates.py}}
\lstinputlisting[language=Python]{../scripts/probe_templates.py}
\subsubsection*{\texttt{scripts/run\_all.py}}
\lstinputlisting[language=Python]{../scripts/run_all.py}
\subsubsection*{\texttt{scripts/run\_fitted\_scenario.py}}
\lstinputlisting[language=Python]{../scripts/run_fitted_scenario.py}
\subsubsection*{\texttt{scripts/run\_pipeline.py}}
\lstinputlisting[language=Python]{../scripts/run_pipeline.py}
\subsubsection*{\texttt{scripts/run\_q1.py}}
\lstinputlisting[language=Python]{../scripts/run_q1.py}
\subsubsection*{\texttt{scripts/run\_q123\_accuracy.py}}
\lstinputlisting[language=Python]{../scripts/run_q123_accuracy.py}
\subsubsection*{\texttt{scripts/run\_q2.py}}
\lstinputlisting[language=Python]{../scripts/run_q2.py}
\subsubsection*{\texttt{scripts/run\_q3.py}}
\lstinputlisting[language=Python]{../scripts/run_q3.py}
\subsubsection*{\texttt{scripts/run\_q4.py}}
\lstinputlisting[language=Python]{../scripts/run_q4.py}
\subsubsection*{\texttt{scripts/selfcheck.py}}
\lstinputlisting[language=Python]{../scripts/selfcheck.py}
\subsubsection*{\texttt{scripts/validate\_accuracy.py}}
\lstinputlisting[language=Python]{../scripts/validate_accuracy.py}
\subsubsection*{\texttt{scripts/validate\_fitted\_scenario.py}}
\lstinputlisting[language=Python]{../scripts/validate_fitted_scenario.py}
\subsubsection*{\texttt{scripts/validate\_latent.py}}
\lstinputlisting[language=Python]{../scripts/validate_latent.py}
\subsubsection*{\texttt{scripts/verify\_accuracy\_outputs.py}}
\lstinputlisting[language=Python]{../scripts/verify_accuracy_outputs.py}
\subsubsection*{\texttt{tests/test\_accuracy.py}}
\lstinputlisting[language=Python]{../tests/test_accuracy.py}
\subsubsection*{\texttt{tests/test\_framework.py}}
\lstinputlisting[language=Python]{../tests/test_framework.py}
\subsubsection*{\texttt{tests/test\_numerics.py}}
\lstinputlisting[language=Python]{../tests/test_numerics.py}
\subsubsection*{\texttt{paper/build\_sources.py}}
\lstinputlisting[language=Python]{../paper/build_sources.py}

